\par\noindent\textbf{(1)} 线化方程为u\_t=−p\_x/ρ0、w\_t=−p\_z/ρ0+b、b\_t=−N²w和u\_x+w\_z=0；N²>0代表密度随高度下降的稳定层结。这里浮力扰动b不是背景磁场，连续性条件把横向与竖向速度耦合。
\par\noindent\textbf{(2)} 取平面波后连续性给ku+mw=0，横向动量给p/ρ0=ωu/k。将浮力b=N²w/(iω)代回竖向动量，消除u和p可得ω²=N²k²/(k²+m²)。此色散具有方向依赖而不是声波式ω∝√(k²+m²)。
\par\noindent\textbf{(3)} 在k>0的正频率支，ω=Nk/√(k²+m²)，群速度为(v\_gx,v\_gz)=N(m²,−km)/(k²+m²)\textasciicircum{}\{3/2\}。点积kv\_gx+mv\_gz=N(km²−km²)/(k²+m²)\textasciicircum{}\{3/2\}=0，故群能流方向与相位波矢垂直。
\par\noindent\textbf{(4)} m→0时ω→N且群速趋零，k→0时ω→0。若背景反向层结使N²<0，则ω虚化，产生对流型增长而非稳定内部波。旋转、有限深度、黏性和非Boussinesq效应均会修改当前结果。
